% =========================================================
% Volume III — Category Theory
% =========================================================
\chapter{Category Theory}

\begin{tcolorbox}[
  colback=gray!6,
  colframe=gray!40,
  arc=2pt,
  left=8pt, right=8pt, top=6pt, bottom=6pt,
  title={\small\textbf{Where You Are in the Journey}},
  fonttitle=\small\bfseries
]
\begin{center}
\small
Algebraic Structures
$\;\to\;$ \textbf{Category Theory} $\;\leftarrow$ you are here
$\;\to\;$ Linear Algebra
$\;\to\;$ Abstract Algebra
$\;\to\;$ Analysis
\end{center}
\end{tcolorbox}

\vspace{0.5em}

\begin{tcolorbox}[
  colback=blue!4,
  colframe=blue!30,
  arc=2pt,
  left=8pt, right=8pt, top=6pt, bottom=6pt,
  title={\small\textbf{Chapter Overview}},
  fonttitle=\small\bfseries
]
\small
Category theory is the abstract framework that unifies all of mathematics
by focusing on \emph{structure-preserving maps} rather than the elements
of individual structures.
A category consists of objects and morphisms (arrows) satisfying
composition and identity laws.

\medskip
\textbf{Why it belongs here:}
Every structure studied in this volume is an example of a category.
Groups with group homomorphisms, rings with ring homomorphisms,
topological spaces with continuous maps, metric spaces with
isometries --- all are categories.
Category theory gives a single language for all of them.

\medskip
\textbf{Connection to the goal problem:}
Stochastic analysis on manifolds (Hsu) uses categorical language implicitly.
Functors formalise ``structure-preserving constructions'';
natural transformations formalise ``coherent families of maps.''
The language will become natural before you need it.

\medskip
\textbf{Primary source:} Mac Lane, \textit{Categories for the Working Mathematician}.
Accessible entry: Awodey, \textit{Category Theory} (Oxford, free draft online).
\end{tcolorbox}

\vspace{0.5em}

\begin{tcolorbox}[
  colback=amber!6,
  colframe=amber!40,
  arc=2pt,
  left=8pt, right=8pt, top=6pt, bottom=6pt,
  title={\small\textbf{Planned Content}},
  fonttitle=\small\bfseries
]
\small
\begin{enumerate}
  \item \textbf{Categories} --- objects, morphisms, composition,
        identity, associativity; concrete examples
  \item \textbf{Functors} --- covariant and contravariant; examples:
        forgetful functor, free functor, hom-functor
  \item \textbf{Natural Transformations} --- definition,
        naturality square; examples from algebra and topology
  \item \textbf{Universal Properties} --- products, coproducts,
        initial and terminal objects; limits and colimits
  \item \textbf{Adjoints} --- adjoint functors, unit and counit,
        the adjoint functor theorem; examples: free/forgetful,
        product/$\Sigma$
  \item \textbf{Examples in This Volume} --- $\mathbf{Grp}$,
        $\mathbf{Ring}$, $\mathbf{Vect}$, $\mathbf{Top}$,
        $\mathbf{Met}$; how each chapter's structure maps here
\end{enumerate}
\end{tcolorbox}

% ---- Notes ----
% =========================================================
% Category Theory — Notes Index
% =========================================================

% (stub — notes to be added as chapter is developed)

